\startenvironment

  %%%%% FONTS %%%%%
  % TODO: I don't understand enough to know why the `use<something>` command is
  % necessary before the two `\define<...>family` commands...
  \usebodyfont[]

  \definefallbackfamily[main][rm][BIZUDPMincho][preset=range:japanese]
  \definefontfamily[main][rm][LinguisticsPro]
  
  % Enable for line breaking across CJK words
  \setscript[nihongo]
  
  % NB: Each family only supports certain sizes; lookup which ones are available.
  \setupbodyfont[main,rm,12pt]
  
  %%%%% TABLES %%%%%
  \definecharacter NUT \utfchar{0x2002}

  % Graphically draw the outer borders, since rule thickness would apply to all sides
  \startuseMPgraphic{MP:thickRight}
    draw rightboundary OverlayBox withpen pensquare scaled 3pt;
    clip currentpicture to OverlayBox;
    setbounds currentpicture to OverlayBox;
  \stopuseMPgraphic
  \startuseMPgraphic{MP:thickTop}
    draw topboundary OverlayBox withpen pensquare scaled 3pt;
    clip currentpicture to OverlayBox;
    setbounds currentpicture to OverlayBox;
  \stopuseMPgraphic
  
  \startuseMPgraphic{MP:upperCorner}
    draw topboundary OverlayBox withpen pensquare scaled 4pt;
    draw rightboundary OverlayBox withpen pensquare scaled 4pt;
    clip currentpicture to OverlayBox;
    setbounds currentpicture to OverlayBox;
  \stopuseMPgraphic
  
  \defineoverlay[OL:thickRight][\useMPgraphic{MP:thickRight}]
  \defineoverlay[OL:thickTop][\useMPgraphic{MP:thickTop}]
  \defineoverlay[OL:upperCorner][\useMPgraphic{MP:upperCorner}]
  
  \setupTABLE[
    direction=reverse,
    rulethickness=0.25pt,
    alignmentcharacter={text->\NUT{}},
    aligncharacter=yes,
    align={middle,lohi},
    offset=0.5em]
  \setupTABLE[row][first][background={OL:thickTop}]
  \setupTABLE[column][first][background={OL:thickRight}]
  \setupTABLE[1][1][background={OL:upperCorner}]
  \setuplegend[before=\startlocalfootnotes, after=\stoplocalfootnotes]
  % \setupcaption[table][location=top, align=flushleft, width=max, minwidth=\textwidth]
  
  % Use symbol-based numbering
  \setupnotation[footnote][numberconversion=set 2]
  
  %%%%% LAYOUT %%%%%
  \setupwhitespace[big]   % Paragraph spacing
  \setupinterlinespace[line=4.0ex]  % Line spacing
  \setuppapersize[letter]
  \setuplayout[
    topspace=0.25in,
    backspace=0.5in,
    edge=0.0in,
    margin=0.0in,
    footer=0.0in,
    footerdistance=0.0in,
    bottom=0.25in,
    % Necessary to make layout fill the entire page.
    width=fit,
    height=fit]
  
  \setupheadertexts % Reset
  \setupheadertexts [chapter] [pagenumber] [pagenumber] [chapter]
  \setuppagenumbering[alternative=doublesided]
    
  % Path is relative to the calling file, not this environment file.
  \setupexternalfigures
    [directory={../media/},
    maxheight=13.0em]
  
  \setupframedtext[align=normal,verytolerant]

  % Keep section headers with some lines of their initial paragraph
  \setuphead[section][before={\testpage[4]}]
  \setuphead[subsection][before={\testpage[2]}]
  \setuphead[section][hang=line]
  
  % Shortcut for temporary language switch inline with main text.
  \define[1]\JPN{\start\language[ja]#1\stop}
  
  % さくら appears a lot, so render her name in a special way.
  \define\Sakura{\dontleavehmode\hbox{\JPN{さくら}}}

  % Custom macro for inserting a glossed example sentence.
  \define[2]\Gloss{
    \startframedtext\JPN{#1}\crlf
      \quote{#2}
    \stopframedtext
  }

  \definedelimitedtext[Grapheme] [left=\utfchar{0x3008}, right=\utfchar{0x3009}]
  \definedelimitedtext[Phoneme] [left=/, right=/]
  \definedelimitedtext[Phonetic] [left=\utfchar{0x005B}, right=\utfchar{0x005D}]

\stopenvironment
